\clearpage
\section{046 -- Banco supercap RTC ridotto da 10 a 8 pezzi}

\begin{tabularx}{\linewidth}{@{}>{\bfseries}l X@{}}
\toprule
Data & 11/11/2026\\
Revisione & non indicata\\
Redatto da & Claude Code\\
Approvazione & Brik81\\
Stato & \tg{tagdone}{MODIFICA ESEGUITA DA BRIK81} \quad \tg{tagcheck}{DA VERIFICARE IN LAYOUT}\\
\bottomrule
\end{tabularx}

\subsection{Componenti, net e fogli coinvolti}
\begin{itemize}
  \item Foglio \textbf{MCU}, sezione RTC. Net \code{RTC\_VBAT}.
  \item Banco di supercondensatori Seiko \code{CPH3225A} (EDLC \SI{11}{\milli\farad}, \SI{3,3}{\volt}, \SI{3,2}{\milli\metre}$\times$\SI{2,5}{\milli\metre}$\times$\SI{0,9}{\milli\metre}, ESR circa \SI{160}{\ohm}) in parallelo sul VBAT di \code{U404} (DS3231MZ).
  \item Carica da \code{3V3} tramite \code{D401} (BAT54WS, anodo su \code{3V3}, blocca la scarica verso \code{3V3}) e \code{R412} (\SI{1}{\kilo\ohm}, limitazione della corrente di carica).
  \item Riferimenti rimossi: \code{C419} e \code{C424}.
\end{itemize}

\subsection{Cosa è cambiato}
\begin{center}
\small
\begin{tabular}{@{}p{3.6cm}p{5.6cm}p{5.6cm}@{}}
\toprule
Aspetto & Prima & Dopo\\
\midrule
Numero di CPH3225A & 10 & 8\\
Capacità totale & $10\times\SI{11}{\milli\farad}=\SI{110}{\milli\farad}$ & $8\times\SI{11}{\milli\farad}=\SI{88}{\milli\farad}$\\
Pin + su \code{RTC\_VBAT} & \code{C415}\ldots\code{C424} & \code{C415}, \code{C416}, \code{C417}, \code{C418}, \code{C420}, \code{C421}, \code{C422}, \code{C423}\\
Altri membri della net & \code{R412.1}, \code{U404.6} & \code{R412.1}, \code{U404.6} (invariati)\\
\bottomrule
\end{tabular}
\end{center}
La modifica è stata eseguita da Brik81 direttamente in KiCad, non dall'agente. Verifica sulla netlist dopo la modifica: la net \code{RTC\_VBAT} contiene esattamente i pezzi elencati sopra. L'agente aveva proposto di togliere \code{C423}/\code{C424}; la scelta dei riferimenti (\code{C419} e \code{C424}) è di Brik81 e il risultato è equivalente.

\subsection{Motivazione}
\textbf{Analisi (verificata su netlist esportata).} Il banco alimenta \emph{solo} il VBAT del DS3231MZ. Non può mantenere in vita l'MCU (STM32L562, \code{U301}, con VBAT collegato direttamente a \code{3V3}) né la FRAM (\code{U401}, FM25V10). L'idea iniziale di \SI{10}{\second} di backup per tutta la scheda è irrealizzabile:
\begin{itemize}
  \item 10 CPH contengono in totale circa \SI{0,5}{\joule};
  \item l'ESR parallela (circa \SI{16}{\ohm}) fa crollare la tensione già a decine di \si{\milli\ampere};
  \item nell'ingombro 3225 / \SI{0,9}{\milli\metre} non esiste un supercap o una batteria a bassa ESR (CPX3225A752D: \SI{7,5}{\milli\farad}, \SI{2,6}{\volt}, \SI{25}{\ohm}, obsoleto; batterie a stato solido TDK CeraCharge da \SI{1,1}{\milli\metre} e Cymbet CBC050 erogano solo da \si{\micro\ampere} a \si{\milli\ampere}).
\end{itemize}
Dopo una riunione Brik81 ha chiarito che il requisito è tenere in vita \emph{solo l'RTC} per \SI{1}{\hour}, al massimo \SI{2}{\hour}.

\textbf{Dimensionamento.}
\begin{itemize}
  \item Datasheet DS3231 (PDF nel progetto, nota 5): \code{IBATT} è una corrente media che include la conversione di temperatura; \SI{3,0}{\micro\ampere} max a VBAT \SI{3,63}{\volt} (DS3231M: \SI{2}{\micro\ampere} tip., \SI{3,0}{\micro\ampere} max). VBAT min \SI{2,3}{\volt}.
  \item Carica a circa \SI{3,1}{\volt} (\code{3V3} meno Vf del BAT54), quindi $\Delta V\approx\SI{0,8}{\volt}$.
  \item Caso peggiore assunto: capacità $-30\,\%$; corrente totale circa \SIrange{4,6}{4,8}{\micro\ampere}, somma di \SI{3}{\micro\ampere} (RTC), circa \SI{1}{\micro\ampere} di leakage inverso del BAT54 a caldo e circa \SI{0,1}{\micro\ampere} per supercap. \textbf{Le ultime due voci sono stime, non dati da datasheet.}
\end{itemize}
\begin{center}
\begin{tabular}{@{}lcc@{}}
\toprule
Numero di CPH3225A & Autonomia caso peggiore & Autonomia tipica\\
\midrule
6 & circa \SI{2,2}{\hour} & \SIrange{6}{7}{\hour}\\
8 & circa \SI{2,8}{\hour} & circa \SI{9}{\hour}\\
\bottomrule
\end{tabular}
\end{center}
Scelti 8 pezzi (Brik81 voleva un numero pari, 6 o 8).

\textbf{Approvazione di Brik81 (testuale):} ``Dai vedo di riuscire a farne 8 al massimo scendo a 6''; e poi ``ho gia rimosso io fregatene''.

\subsection{Impatti e criticità}
\begin{notebox}
\textbf{Layout.} Se in layout gli 8 pezzi non entrano si scende a 6 (si tolgono altri 2): con 6 pezzi l'autonomia nel caso peggiore è circa \SI{2,2}{\hour}, appena sopra il massimo richiesto di \SI{2}{\hour}.
\end{notebox}
\begin{notebox}
\textbf{Stime.} I valori del leakage inverso del BAT54 a caldo e della corrente per supercap del CPH3225A sono stime, non letture da datasheet. L'autonomia reale non è ancora nota.
\end{notebox}
Il backup della scheda intera (MCU, FRAM) resta non coperto dal banco, come già deciso nelle entry 044 e 045.

\subsection{Verifiche da fare}
\begin{itemize}
  \item Controllare in layout che gli 8 CPH3225A entrino; in caso contrario scendere a 6.
  \item Misurare l'autonomia reale dell'RTC sulla prima versione hardware (Brik81 non può testare prima).
  \item Verificare se il firmware salva comunque dati in FRAM per altri motivi (non richiede backup hardware).
\end{itemize}
